\section{Multi-GPU Oceananigans on Alliance Clusters}
\label{app:oceananigans-alliance}

This appendix gets Oceananigans.jl running on the GPUs of the Alliance clusters, from a first test
job to multi-node runs with CUDA-aware MPI. Everything that can change, such as module versions
and MPI settings, lives in a tested repository, so the steps here stay short:

\begin{center}
\url{https://github.com/francispoulin/Oceananigans-DRAC}
\end{center}

It was tested on Nibi and Fir in October 2026 (Table~\ref{tab:alliance-settings}). New to the
clusters? Read Appendix~\ref{app:alliance} first, and Appendix~\ref{app:ssh} for logging in with keys.

\subsection{Before You Start}

\begin{itemize}
\item An Alliance account, multi-factor authentication, and your supervisor's account name
(written \texttt{def-yourpi} below).
\item A \emph{minimal} \texttt{\textasciitilde/.bashrc}: no \texttt{module load} lines and no
\texttt{LD\_LIBRARY\_PATH}. The environment script loads what it needs. Also check that
\texttt{\textasciitilde/.modulerc} does not exist or pin an old \texttt{StdEnv}.
\item Install packages on a \emph{login node}: login nodes have internet but no GPUs, so
\texttt{nvidia-smi} failing there is normal.
\end{itemize}

\subsection{Get the Code}

\begin{lstlisting}
cd $SCRATCH
git clone https://github.com/francispoulin/Oceananigans-DRAC.git
# or, with a GitHub SSH key (Appendix C):
# git clone git@github.com:francispoulin/Oceananigans-DRAC.git
cd Oceananigans-DRAC
\end{lstlisting}

\subsection{The Environment}

\begin{lstlisting}
source env/drac.sh
\end{lstlisting}

This loads the modules (\texttt{StdEnv/2023}, \texttt{gcc}, \texttt{openmpi}, \texttt{cuda},
\texttt{julia}) and sets the Julia \emph{depot}, the folder where packages and compiled code are
kept. The default is \texttt{\$SCRATCH/julia\_depot}, because home directories have small quotas.
To use another folder, set \texttt{JULIA\_DEPOT\_PATH} \emph{before} sourcing the script. Source it in
every new shell; the repository's job scripts do it for you. The line it prints shows the cluster,
the repository and the depot, which is the first thing to check when something looks wrong.

\subsection{One-Time Setup}

\begin{lstlisting}
bash setup/setup.sh
\end{lstlisting}

This takes 20--30 minutes. It installs the tested package versions and makes both MPI.jl and
Julia's own OpenMPI use the cluster's CUDA-aware OpenMPI. In its output, the paths for
\texttt{libmpi} must be under \texttt{/cvmfs}, not in your depot, and there should be no errors.

\subsection{Test Jobs}

\begin{lstlisting}
sbatch --account=def-yourpi jobs/hello_cpu.sh
sbatch --account=def-yourpi jobs/hello_gpu.sh
sbatch --account=def-yourpi --nodes=1 --ntasks-per-node=2 jobs/checks.sh
\end{lstlisting}

Each job writes a \texttt{.out} file. The hello jobs end with \texttt{hello\_oceananigans: OK}.
\texttt{checks.sh} prints PASS or FAIL for each check and ends with \texttt{All checks passed}.
It tests MPI, CUDA-aware communication between GPUs, a distributed run against a single-GPU run,
and NetCDF and JLD2 output. To test a full node or several nodes, use the settings in
Table~\ref{tab:alliance-settings}, for example:

\begin{lstlisting}
sbatch --account=def-yourpi --nodes=2 --ntasks-per-node=4 --mem=0 jobs/checks.sh
\end{lstlisting}

\subsection{Your Own Scripts}

Copy \texttt{jobs/checks.sh} and replace its loop of checks with your command:

\begin{lstlisting}
srun julia --project=. my_script.jl
\end{lstlisting}

Run one MPI rank per GPU (\texttt{--ntasks-per-node} equal to the GPUs you want per node, with
\texttt{--gpus-per-task=h100:1}). In the Julia script, create the architecture with
\texttt{arch = Distributed(GPU())}. The scripts in \texttt{checks/} are working examples.

\begin{table}[h]
\centering\small
\begin{tabular}{@{}lcll@{}}
\toprule
\textbf{Cluster} & \textbf{GPUs/node} & \textbf{Full node} & \textbf{Status} \\
\midrule
Nibi     & 8 & \texttt{--nodes=1 --ntasks-per-node=8 --mem=0} & tested \\
Fir      & 4 & \texttt{--nodes=1 --ntasks-per-node=4 --mem=0} & tested \\
Rorqual  & 4 & \texttt{--nodes=1 --ntasks-per-node=4 --mem=0} & not yet tested \\
Trillium & 4 & see its documentation                          & not yet tested \\
\bottomrule
\end{tabular}
\caption{Job settings by cluster (October 2026). For several nodes, increase \texttt{--nodes}.}
\label{tab:alliance-settings}
\end{table}

\subsection{Common Problems}

\paragraph{Julia cannot find my packages, or the wrong Julia starts.}
Every new shell needs \texttt{source env/drac.sh}. Check with:

\begin{lstlisting}
which julia; julia --version; echo $JULIA_DEPOT_PATH
\end{lstlisting}

You should see Julia 1.10.10 and your depot folder. Another version usually means that
\texttt{module load} lines in \texttt{\textasciitilde/.bashrc}, or a \texttt{\textasciitilde/.modulerc}
that sets an old default \texttt{StdEnv}, are interfering.

\paragraph{Precompilation fails with ``could not spawn''.}
Two usual causes. (1) An old \texttt{LD\_LIBRARY\_PATH} in \texttt{.bashrc}: \texttt{echo
\$LD\_LIBRARY\_PATH} should print nothing. (2) Login-node limits: retry with fewer parallel
tasks, or precompile inside an interactive allocation.

\begin{lstlisting}
JULIA_NUM_PRECOMPILE_TASKS=2 julia --project=. -e 'using Pkg; Pkg.precompile()'
salloc --account=def-yourpi --cpus-per-task=4 --mem=16G --time=1:00:00
\end{lstlisting}

\paragraph{A GPU job fails with ``Out of GPU memory'' on a tiny problem.}
The GPU may be occupied by another process or in a bad state. Look at the line printed by
\texttt{hello\_gpu.sh} for the memory in use at the start. Rerun; if it fails again on the same
node, send the node name and job ID to Alliance support.

\paragraph{A job stays pending.}
\begin{lstlisting}
squeue -u $USER --start
\end{lstlisting}
shows the estimated start. Whole-node requests wait longer than small ones.

\paragraph{My packages disappeared.}
Scratch is temporary storage, and unused files are eventually deleted. If your depot is gone,
source the environment and rerun \texttt{bash setup/setup.sh}.
